% GENERATED FILE. Edit canonical lessons, then run npm run generate:books.
% canonical-source-hash: fnv1a64:ed0d0af0a6c62da6
% canonical-lessons: PA-W01-na, PA-W01-ii-matra, PA-W01-ma, PA-W01-sa, PA-W01-ta, PA-W01-e-matra, PA-W01-namaste-read

\chapter{From the Mouth of the Guru, One Piece at a Time}
\label{ch:pa-script-first}
\hlchaptermodality{14}

\begin{chapteropening}
\textbf{By the end of this chapter:} \emph{I can add six Gurmukhi pieces to the three learned with \pa{ਹਾਂ}, then read \pa{ਨਹੀਂ} and \pa{ਨਮਸਤੇ} by assembling only known pieces.}

The greeting assembled from five pieces the reader can each write from memory, introducing nothing new, and naming plainly the one piece of machinery still to come.
\end{chapteropening}

\section[na]{\pa{ਨ} — the consonant na}

\label{lesson:PA-W01-na}



\begin{quote}
Write the piece from the previous lesson before adding another.
\end{quote}

\subsection*{Script}
\begin{quote}
\pa{ਨ}
\end{quote}

\textbf{na} — the \emph{n} of English \emph{net}, with the inherent \textbf{a} riding along.

No new machinery. Every rule you already hold applies to it unchanged, and that is the point of learning the system rather than the words: \textbf{new shapes are coming; new rules are not.}

\subsection*{Your turn}
\begin{itemize}
\raggedright
  \item \emph{Write it:} \pa{ਨ}
  \item \emph{Read it:} every piece this chapter has given you so far, in order
\end{itemize}

\begin{tcolorbox}[breakable,colback=teal!4,colframe=teal!35!black,title={Before you move on}]
What vowel is inside \textbf{\pa{ਨ}} when nothing is attached? (\textbf{a}.) How many new rules did this lesson add? (\textbf{None}.)

Source: \href{https://www.unicode.org/charts/PDF/U0A00.pdf}{Unicode Gurmukhi chart}.
\end{tcolorbox}
\section[-ī]{\pa{ੀ} — a second mātrā}

\label{lesson:PA-W01-ii-matra}



\begin{quote}
Write the piece from the previous lesson before adding another.
\end{quote}

\subsection*{Script}
\begin{quote}
\pa{ੀ}
\end{quote}

\textbf{ī}, the long \emph{ee} of English \emph{see}. Like the first mātrā it stands to the right, but it carries a curve that rises and \textbf{hooks back over the head-line}.

Now build with everything you have:

\begin{quote}
\textbf{\pa{ਨ}} + \textbf{\pa{ਹ}} + \textbf{\pa{ੀ}} + \textbf{\pa{ਂ}} $\to$ \textbf{\pa{ਨਹੀਂ}}
\end{quote}

\textbf{na-hīṁ} — \textbf{no}. Two consonants, one mātrā, one nasal dot.

Both answers to a yes-or-no question are now readable:

\begin{quote}
\textbf{\pa{ਹਾਂ}}   \textbf{\pa{ਨਹੀਂ}}
\end{quote}

That is the arithmetic of an abugida: consonants and mātrās \textbf{multiply}, so a handful of pieces reaches a large number of syllables very quickly.

\subsection*{Your turn}
\begin{itemize}
\raggedright
  \item \emph{Write it:} \pa{ੀ}
  \item \emph{Read it:} every piece this chapter has given you so far, in order
\end{itemize}

\begin{tcolorbox}[breakable,colback=teal!4,colframe=teal!35!black,title={Before you move on}]
What is \textbf{\pa{ਨਹੀਂ}}? (\textbf{No}.) How many pieces is it built from? (\textbf{Four}.) And what is the dot doing at the end? (\textbf{Nasalising} the vowel.)

Source: \href{https://www.unicode.org/charts/PDF/U0A00.pdf}{Unicode Gurmukhi chart}.
\end{tcolorbox}
\section[ma]{\pa{ਮ} — the consonant ma}

\label{lesson:PA-W01-ma}



\begin{quote}
Write the piece from the previous lesson before adding another.
\end{quote}

\subsection*{Script}
\begin{quote}
\pa{ਮ}
\end{quote}

\textbf{ma} — the \emph{m} of English \emph{met}, plus the inherent vowel.

Three consonants. Say them in the order you met them, each with the vowel nobody wrote:

\begin{quote}
\textbf{\pa{ਹ}}   \textbf{\pa{ਨ}}   \textbf{\pa{ਮ}}
\end{quote}

\emph{ha, na, ma.}

\subsection*{Your turn}
\begin{itemize}
\raggedright
  \item \emph{Write it:} \pa{ਮ}
  \item \emph{Read it:} every piece this chapter has given you so far, in order
\end{itemize}

\begin{tcolorbox}[breakable,colback=teal!4,colframe=teal!35!black,title={Before you move on}]
Say the three consonants you hold, each with its inherent vowel. (\emph{\textbf{ha, na, ma}}.)

Source: \href{https://www.unicode.org/charts/PDF/U0A00.pdf}{Unicode Gurmukhi chart}.
\end{tcolorbox}
\section[sa]{\pa{ਸ} — the consonant sa}

\label{lesson:PA-W01-sa}



\begin{quote}
Write the piece from the previous lesson before adding another.
\end{quote}

\subsection*{Script}
\begin{quote}
\pa{ਸ}
\end{quote}

\textbf{sa} — the \emph{s} of English \emph{sit}, plus the inherent vowel.

Four consonants, two mātrās and a nasal dot. One shape and one mātrā left before the greeting is readable.

\subsection*{Your turn}
\begin{itemize}
\raggedright
  \item \emph{Write it:} \pa{ਸ}
  \item \emph{Read it:} every piece this chapter has given you so far, in order
\end{itemize}

\begin{tcolorbox}[breakable,colback=teal!4,colframe=teal!35!black,title={Before you move on}]
What is the inherent vowel in \textbf{\pa{ਸ}}? (\textbf{a} — \emph{sa}.)

Source: \href{https://www.unicode.org/charts/PDF/U0A00.pdf}{Unicode Gurmukhi chart}.
\end{tcolorbox}
\section[ta]{\pa{ਤ} — the consonant ta}

\label{lesson:PA-W01-ta}



\begin{quote}
Write the piece from the previous lesson before adding another.
\end{quote}

\subsection*{Script}
\begin{quote}
\pa{ਤ}
\end{quote}

\textbf{ta} — and it is not quite the English consonant.

English \emph{t} is made with the tongue on the \textbf{ridge behind} the teeth. This one is \textbf{dental}: the tongue touches the \textbf{back of the upper teeth themselves}, further forward. It sounds softer to an English ear, and getting it right is one of the quickest ways to stop sounding foreign.

The script has a \textbf{second} t made further back, and the two are different letters in different words. This is the forward one.

\subsection*{Your turn}
\begin{itemize}
\raggedright
  \item \emph{Write it:} \pa{ਤ}
  \item \emph{Read it:} every piece this chapter has given you so far, in order
\end{itemize}

\begin{tcolorbox}[breakable,colback=teal!4,colframe=teal!35!black,title={Before you move on}]
Where does the tongue go for \textbf{\pa{ਤ}}? (Against the \textbf{back of the upper teeth} — dental.) Is there another t? (\textbf{Yes}, one made further back.)

Source: \href{https://www.unicode.org/charts/PDF/U0A00.pdf}{Unicode Gurmukhi chart}.
\end{tcolorbox}
\section[-e]{\pa{ੇ} — a mātrā that sits ABOVE the letter, not beside it}

\label{lesson:PA-W01-e-matra}



\begin{quote}
Write the piece from the previous lesson before adding another.
\end{quote}

\subsection*{Script}
\begin{quote}
\pa{ੇ}
\end{quote}

\textbf{e}, the vowel of English \emph{they}. And unlike both mātrās you already hold, which stand to the right, this one sits \textbf{above} the letter, riding on the head-line:

\begin{quote}
\textbf{\pa{ਤ}} + \textbf{\pa{ੇ}} $\to$ \textbf{\pa{ਤੇ}}
\end{quote}

\emph{te.} Still one beat.

\textbf{Where a mātrā sits is part of its identity}, not a style choice. In this family of scripts a vowel sign can go right, left, above, below, or on both sides at once, and a reader learns to look \textbf{around} a consonant rather than only after it.

\subsection*{Your turn}
\begin{itemize}
\raggedright
  \item \emph{Write it:} \pa{ੇ}
  \item \emph{Read it:} every piece this chapter has given you so far, in order
\end{itemize}

\begin{tcolorbox}[breakable,colback=teal!4,colframe=teal!35!black,title={Before you move on}]
Where does this mātrā sit? (\textbf{Above} the consonant.) Does it add a beat? (\textbf{No}.)

Source: \href{https://www.unicode.org/charts/PDF/U0A00.pdf}{Unicode Gurmukhi chart}.
\end{tcolorbox}
\section[namaste]{\pa{ਨਮਸਤੇ} — five pieces, all of them already yours}

\label{lesson:PA-W01-namaste-read}



\begin{quote}
Write the piece from the previous lesson before adding another.
\end{quote}

\subsection*{Script}
Nothing new is introduced here. Every piece below is one you can write.

\begin{quote}
\textbf{\pa{ਨ}} + \textbf{\pa{ਮ}} + \textbf{\pa{ਸ}} + \textbf{\pa{ਤ}} + \textbf{\pa{ੇ}} $\to$ \textbf{\pa{ਨਮਸਤੇ}}
\end{quote}

Read it the way you built it:

\begin{itemize}
\raggedright
  \item \textbf{\pa{ਨ}} — \emph{na}, inherent vowel
  \item \textbf{\pa{ਮ}} — \emph{ma}, inherent vowel
  \item \textbf{\pa{ਸ}} — \emph{sa}, inherent vowel
  \item \textbf{\pa{ਤੇ}} — \emph{ta} with the mātrā above, making \emph{te}
\end{itemize}

\textbf{na-ma-ste.} Four consonants in a row, three of them carrying the vowel nobody wrote, and one mark on the last. The greeting from your first lesson, decoded rather than recalled.

\textbf{One piece of machinery is still missing, and it is worth naming rather than hiding.} Some words press two consonants together with no vowel between them, and this script has a way to write that — a mark that deletes the inherent vowel and fuses the pair into one shape. You have not needed it yet. The longer Sikh greeting you already say contains one, and it is the next thing this book will teach you.

So: \textbf{four consonants, three mātrās and a nasal dot}, and with them you can read \emph{yes}, \emph{no} and the greeting.

\subsection*{Your turn}
\begin{itemize}
\raggedright
  \item \emph{Write it:} \pa{ਨਮਸਤੇ}
  \item \emph{Read it:} every piece this chapter has given you so far, in order
\end{itemize}

\begin{tcolorbox}[breakable,colback=teal!4,colframe=teal!35!black,title={Before you move on}]
How many new pieces were in this lesson? (\textbf{None}.) What carries the vowel in \textbf{\pa{ਨ}}, \textbf{\pa{ਮ}} and \textbf{\pa{ਸ}}? (\textbf{Nothing} — it is the inherent vowel.) And what is the one piece of machinery still to come? (The mark that \textbf{deletes} the inherent vowel and \textbf{joins} two consonants.)

Source: \href{https://www.unicode.org/charts/PDF/U0A00.pdf}{Unicode Gurmukhi chart}.
\end{tcolorbox}
